\documentclass[11pt]{article}

\usepackage[a4paper,margin=2.4cm]{geometry}
\usepackage{amsmath}
\usepackage{booktabs}
\usepackage{enumitem}
\usepackage{graphicx}
\usepackage[hidelinks]{hyperref}
\usepackage{siunitx}

\IfFileExists{docs/figures/xsection_fit/channels_compact.pdf}{
  \newcommand{\figuredir}{docs/figures/xsection_fit}
}{
  \newcommand{\figuredir}{figures/xsection_fit}
}

\newcommand{\file}[1]{\texttt{#1}}

\title{Sampling positron-emitter production\\from fitted cross sections}
\author{J.J. G\'omez-Cadenas}
\date{\today}

\begin{document}
\maketitle

\begin{abstract}
Five fitted excitation functions, with correlated uncertainty ensembles,
describe the production of $^{11}$C, $^{13}$N and $^{15}$O by protons on the
nuclei of tissue.  This note explains, step by step, how those curves replace
the internal nuclear model of the transport code as the source of
positron-emitter production, while the transport itself is unchanged.  The
method is a per-step sampling: at every proton step, the fitted cross section
decides whether an emitter is born there.  The deliverable is the paired
comparison of the distal activity edge computed both ways,
$\Delta = R_{50}^{\mathrm{data}} - R_{50}^{\mathrm{G4}}$, with identical
proton histories, so the difference measures the production model alone.
\end{abstract}

\section{Starting point: five fitted excitation functions}

The production of positron emitters by a therapeutic proton beam in tissue is
dominated by five reaction channels,
\begin{equation}
\begin{split}
  &^{12}\mathrm{C}(p,x)^{11}\mathrm{C},\qquad
   ^{16}\mathrm{O}(p,x)^{15}\mathrm{O},\qquad
   ^{16}\mathrm{O}(p,x)^{11}\mathrm{C},\\
  &^{14}\mathrm{N}(p,x)^{13}\mathrm{N},\qquad
   ^{16}\mathrm{O}(p,x)^{13}\mathrm{N}.
\end{split}
\label{eq:channels}
\end{equation}
For each channel $c$ we possess a fitted excitation function
$\sigma_c(E)$ --- the production cross section as a function of proton kinetic
energy --- built from curated experimental data, together with an ensemble of
1000 correlated replica curves that represents its uncertainty.  The
construction of these curves (data selection, fitting, replica generation) is
documented in \file{docs/xsection\_fit.tex}.  Figure~\ref{fig:curves} shows
the five nominal curves with their 16--84\% uncertainty bands and the
measurements they are built from; Fig.~\ref{fig:anatomy} shows the complete
diagnostic anatomy of one fit.

\begin{figure}[htbp]
\centering
\includegraphics[width=\textwidth]{\figuredir/channels_compact.pdf}
\caption{The five fitted excitation functions.  Grey points: accepted
experimental measurements.  Blue band: 16--84\% interval of the replica
ensemble.  Black curve: the nominal fit (pointwise replica median).}
\label{fig:curves}
\end{figure}

\begin{figure}[htbp]
\centering
\includegraphics[width=\textwidth]{\figuredir/p_C12_x_C11.pdf}
\caption{Anatomy of one fit, $^{12}$C(p,x)$^{11}$C: the fitted data and
band, nine representative replicas, the replica-distance distribution, and
the evaluated libraries as external comparisons.  The distance $D$ is the
covariance-weighted mean-square separation between a replica curve and the
nominal, evaluated at the fitted data energies
(\file{docs/xsection\_fit.tex}); the nine displayed replicas sit at fixed
percentiles of its distribution.}
\label{fig:anatomy}
\end{figure}

Each channel opens at its physical threshold, the lowest proton energy at
which the reaction is energetically possible
(Table~\ref{tab:thresholds}).  The curves are zero at and below threshold and
are fitted from threshold to \SI{150}{MeV}, which covers the treatment field
with margin.

\begin{table}[htbp]
\centering
\caption{Adopted laboratory thresholds (AME2020 masses).}
\label{tab:thresholds}
\begin{tabular}{llS[table-format=2.2]}
\toprule
Channel & Lowest reaction & {Threshold (MeV)} \\
\midrule
$^{12}$C(p,x)$^{11}$C  & $^{12}$C(p,d)$^{11}$C            & 17.90 \\
$^{16}$O(p,x)$^{15}$O  & $^{16}$O(p,d)$^{15}$O            & 14.29 \\
$^{16}$O(p,x)$^{11}$C  & $^{16}$O(p,d+$\alpha$)$^{11}$C   & 25.17 \\
$^{14}$N(p,x)$^{13}$N  & $^{14}$N(p,d)$^{13}$N            &  8.90 \\
$^{16}$O(p,x)$^{13}$N  & $^{16}$O(p,$\alpha$)$^{13}$N     &  5.55 \\
\bottomrule
\end{tabular}
\end{table}

\section{What we want to compute}

The clinical observable of PET range verification is the distal edge of the
activity profile: the depth at which the emitter production falls to half of
its maximum, $R_{50}$.  The transport code (Geant4) predicts this edge using
its own internal nuclear model for the reactions in
Eq.~\eqref{eq:channels}.  The fitted curves of Fig.~\ref{fig:curves} are the
experimental knowledge of the same reactions.  The question this note
answers is:
\begin{quote}
\emph{by how much does the predicted distal edge move when the production
cross sections of the internal model are replaced by the fitted experimental
curves, everything else being identical?}
\end{quote}
The answer is one number per isotope (and one for the combined profile),
\begin{equation}
  \Delta_j = R_{50,j}^{\mathrm{data}} - R_{50,j}^{\mathrm{G4}},
  \label{eq:delta}
\end{equation}
where ``data'' denotes production sampled from the fitted curves and ``G4''
denotes production from the internal model, both evaluated on the same proton
histories.

\section{Production probability of one proton step}

The transport code advances each proton (primary or secondary) in short
steps.  Consider one step of length $\ell$ (cm) through a material that
contains $n_t$ nuclei per cm$^3$ of target isotope $t$, taken by a proton of
kinetic energy $E$ and statistical weight $a$.  The expected number of
reactions through channel $c$ on this step is the number of target nuclei per
unit area seen by the proton, multiplied by the cross section:
\begin{equation}
  \mu_{c} = a\, n_{t(c)}\,\ell\,\sigma_c(E),
  \label{eq:mu}
\end{equation}
with $\sigma_c$ in cm$^2$ ($1~\mathrm{mb} = 10^{-27}\,\mathrm{cm}^2$).  For a
therapeutic beam $\mu_c$ is of order $10^{-6}$ per step: the expected count
is far below one, so $\mu_c$ is also, to excellent accuracy, the
\emph{probability} that one reaction happens on the step,
\begin{equation}
  P(\text{one reaction on the step}) = 1 - e^{-\mu_c} \simeq \mu_c .
  \label{eq:bernoulli}
\end{equation}

This gives the sampling rule.  At every proton step, for every channel open
at that energy, draw a uniform random number $u \in (0,1)$; if
$u < \mu_c$, one emitter of the channel's isotope is created at a point
chosen uniformly along the step.  The proton itself continues exactly as
before --- the sampled reaction changes nothing in its transport.  The
created emitter then follows its natural fate: the isotope decays with its
physical half-life, the positron slows down in the surrounding material and
annihilates, and the production and annihilation points are recorded.

Summing Eq.~\eqref{eq:mu} over all steps of all protons gives the expected
production of channel $c$,
\begin{equation}
  N_c = \sum_{s} a_s\, n_{t(c),s}\,\ell_s\,\sigma_c(E_s),
  \label{eq:total}
\end{equation}
and the sampled emitters are a Poisson realization of this expectation,
distributed in space exactly as production demands: where the proton fluence
is high, where the target nuclei are abundant, and where the cross section at
the local proton energy is large.

\section{Two tallies on the same protons}

The comparison runs both production models on \emph{identical} proton
histories.  During the same transport:
\begin{enumerate}[nosep]
  \item the \textbf{data-driven tally} samples emitters with
        Eq.~\eqref{eq:bernoulli} from the fitted curves, as described above;
  \item the \textbf{native tally} counts the residual nuclei that the
        transport code's internal nuclear model produces --- the same five
        channels, recorded with production position and identity.
\end{enumerate}
Because both tallies watch the same protons step by step, the statistical
noise of the transport cancels in the comparison: $\Delta_j$ in
Eq.~\eqref{eq:delta} measures the production model, not the luck of the run.

Two rules keep the calculation consistent.  First, the residual nuclei of
the five channels created by the internal model are counted but \emph{not}
turned into emitters --- production of these isotopes is claimed once, by the
fitted curves.  Second, the proton transport is left untouched by the
sampling.  This second rule deserves a closer look.

As the beam penetrates, protons are removed from it by nuclear collisions:
a proton that undergoes an inelastic collision with a nucleus scatters away
and leaves the beam, so the proton fluence decreases with depth --- roughly
one percent of the protons per centimetre of tissue.  The removal rate is
set by the \emph{total} inelastic cross section, the probability of a
nuclear collision with any outcome whatsoever, several hundred millibarns.
The five channels of Eq.~\eqref{eq:channels} are particular outcomes of
such collisions; together they amount to tens of millibarns, the fraction
\begin{equation}
  f(E)=\frac{\sum_c\sigma_c(E)}{\sigma_{\mathrm{inel}}(E)}
  \label{eq:fraction}
\end{equation}
of the total, about ten percent.

Replacing the five partial cross sections while the transport keeps the
total fixed introduces a first-order inconsistency: a channel whose fitted
cross section exceeds the internal model's should, strictly, remove slightly
more protons from the beam.  The size of the inconsistency is set by the
\emph{change} of the partial cross sections --- a few millibarns against
several hundred --- so the removal rate is altered by well below one
percent, and most of that cancels in the shape of the profile that
determines $R_{50}$.  The fraction $f(E)$ is the conservative ceiling of
the effect, and it is tracked as a diagnostic over the energies and
materials sampled by the treatment field.

\section{From production to the distal edge}

Binning the emitters of isotope $j$ in depth $z$ along the beam axis gives
the production profile $P_j(z)$; the sum over isotopes gives the combined
profile.  The distal edge is the depth beyond the profile maximum at which
the profile crosses half of that maximum.  With bin centres $z_k$ and
contents $P_k$, let $k^\ast$ be the last bin with
$P_{k^\ast} \ge \tfrac{1}{2}P_{\max}$; linear interpolation to the next bin
gives
\begin{equation}
  R_{50} = z_{k^\ast}
  + \left(\tfrac{1}{2}P_{\max} - P_{k^\ast}\right)
    \frac{z_{k^\ast+1}-z_{k^\ast}}{P_{k^\ast+1}-P_{k^\ast}} .
  \label{eq:r50}
\end{equation}
Absolute yields are reported per simulated run, per primary proton, and per
gray, using the proton count and the target dose recorded with the run.

\section{Deliverable}

The first sampling run produces, for the standard treatment field:
\begin{enumerate}[nosep]
  \item the production profiles $P_j(z)$ from both tallies, per isotope and
        combined, on the same depth axis;
  \item the per-isotope yields per gray from both tallies;
  \item the paired edge displacements $\Delta_j$ of Eq.~\eqref{eq:delta}.
\end{enumerate}
This section receives those figures and numbers when the first calculation
runs.

The same sampling machinery extends directly to the uncertainty band: each
of the 1000 replica curves defines its own $\sigma_c(E)$, and the spread of
the resulting edges gives the cross-section uncertainty of $R_{50}$.  That
propagation, and its efficient implementation, are specified in
\file{docs/xsections\_plan.tex} and \file{docs/shared\_plan.tex}.

\end{document}
